\documentclass[11pt,aspectratio=169,handout]{beamer} 
\usetheme{Antibes} % si possono modificare in tutti i modi possibili
\usecolortheme{whale}
\usefonttheme[onlymath]{serif}

\definecolor{myTeal}{RGB}{31 122 140}
\definecolor{myMutedTeal}{RGB}{1 22 70}
\definecolor{myBlue}{RGB}{0 123 167}
\definecolor{myLavender}{RGB}{118 108 127}
\definecolor{myApricot}{RGB}{249 219 189}
%\definecolor{myGreen}{RGB}{34 139 34}
%\definecolor{myGreen}{RGB}{0 150 44}
%\definecolor{myGreen}{RGB}{37 110 75}
\definecolor{myGreen}{RGB}{19 96 55}
%\definecolor{myGreen}{RGB}{43 129 88}

\setbeamercolor{structure}{fg=myGreen} % bullet points
\setbeamercolor{palette primary}{bg=myGreen, fg=white} % headers colour
\setbeamercolor{titlelike}{fg=white} % titles colour 

\setbeamertemplate{navigation symbols}{} % nasconde i pulsanti di navigazione
\setbeamertemplate{headline}{} % navigazione in cima

\usepackage{appendixnumberbeamer}
\usepackage{mathtools}
\usepackage{lipsum}

% per l'handout
% \usepackage{pgfpages}
% \pgfpagesuselayout{4 on 1}[a4paper, border shrink = 5mm]

\usepackage[utf8]{inputenc}
\usepackage[italian]{babel}
\usepackage{caption}

\usepackage{chronosys}
\setupchronoevent{textwidth=3cm, textstyle=\centering}

\newtheorem{teo}[theorem]{Teorema}
\newtheorem{de}[theorem]{Definizione}
\newtheorem{no}[theorem]{Nota}
\setbeamercolor{block title}{bg=myGreen,fg=white} 

\newcommand{\COMMENT}[1]{} %comment block
\usepackage{amsmath,physics}

\title{\textbf{Thesis Title}}
\author{Name Surname \\ Relatore/trice: Cool Guy}
\date{67 November 4200}

\begin{document}

\section{Intro}

\begin{frame}[plain] %Frames are slides
    \centering
    \includegraphics[width=2cm]{logounito.png}\par
    \maketitle
\end{frame}
%------------------------------------------------------------------------------
\setbeamertemplate{footline}{ % this sets the headers & footers
    \hspace{0.5cm}
    \includegraphics[width=0.7cm]{logounito.png}
    \hfill%
    \usebeamercolor[fg]{page number in head/foot}%
    \usebeamerfont{page number in head/foot}%
    \insertframenumber\,/\,\inserttotalframenumber\kern1em\vskip2pt% pag/tot
}
%------------------------------------------------------------------------------
\begin{frame}{Introduction}
    \begin{columns}[t]
        \column{.50\textwidth}        
        \phantom{miao}\\
        Sample text Sample text. Sample text. Sample text. Sample text. Sample text. 
        Sample text Sample text. Sample text. Sample text. Sample text. Sample text. 
        Sample text Sample text. Sample text. Sample text. Sample text. Sample text. 
        Sample text Sample text. Sample text. Sample text. Sample text. Sample text. 
        Sample text Sample text. Sample text. Sample text. Sample text. Sample text. 
        Sample text Sample text. Sample text. Sample text. Sample text. Sample text. 
        Sample text Sample text. Sample text. 
        \column{.45\textwidth}
        \begin{block}{\textbf{Indice}} 
            \begin{enumerate}
                \item Algebra di Supersimmetria.
                \item Hamiltoniane \textit{partner}.
                \item \textit{Shape invariance}.
                \item Metodo SWKB.
            \end{enumerate}    
        \end{block}
    \end{columns}
\end{frame}
%------------------------------------------------------------------------------
\begin{frame}{Indice}
    \begin{columns}       
        \column{.55\textwidth}
       \begin{enumerate}
           \item Algebra di Supersimmetria.
               \vspace{0.3cm}
           \item Hamiltoniane \textit{partner} supersimmetriche:
               \begin{itemize}
                   \item Rapporto tra autostati e tra autovalori;
                  \item Spettri discreti;
                  \item Supersimmetria rotta e non rotta;
                  \item Esempio: la buca infinita di potenziale.
                \end{itemize}
        \end{enumerate}
        %\vspace{0.3cm}
        \column{.45\textwidth}
        \begin{enumerate}
            \setcounter{enumi}{2} 
            \item \textit{Shape invariance}:
                \begin{itemize}
                    \item Operatori scala generalizzati;
                    \item Esempio: l'oscillatore armonico isotropo.
                \end{itemize}
                \vspace{0.3cm}
            \item Metodo SWKB:
                \begin{itemize}
                    \item SWKB del primo ordine;
                    \item SWKB per potenziali \textit{shape invariant}.
                \end{itemize}
        \end{enumerate}
    \end{columns}
\end{frame}

%------------------------------------------------------------------------------
\appendix

\setbeamertemplate{footline}{
    \hspace{0.5cm}
    \includegraphics[width=0.7cm]{logounito.png}
    \hfill%
}

\section*{Thank you!}
%------------------------------------------------------------------------------
\begin{frame}[plain]
    \centering \Large
    \vspace{1.5cm}

    Grazie per l'attenzione!

    \vspace{1cm}

    \includegraphics[width=2cm]{logounito.png}
\end{frame}
%------------------------------------------------------------------------------
        \begin{frame}[plain]
            \centering
        \includegraphics[width=2cm]{logounito.png}
        \end{frame}
%------------------------------------------------------------------------------
        

%------------------------------------------------------------------------------
%-------------------------------BACKUP-----------------------------------------
%-------------------------------SLIDES-----------------------------------------
%------------------------------------------------------------------------------
\setbeamertemplate{footline}{ % this sets the headers & footers without number
    \hspace{0.5cm}
    \includegraphics[width=0.7cm]{logounito.png}
    \hfill%
    \usebeamercolor[fg]{page number in head/foot}%
    \usebeamerfont{page number in head/foot}%
    %\insertframenumber\,/\,\inserttotalframenumber\kern1em
    \vskip2pt% pag/tot
}

\begin{frame}{Esempio: la buca infinita di potenziale/1}
    Si consideri la buca infinita di potenziale, con un \textit{offset} energetico
    tale da rendere lo stato fondamentale a energia nulla:
    $$V_1(x)=\begin{dcases}
        -\frac{\pi^2\hbar^2}{2mL^2}&\text{se } 0\leq x\leq L\\
        \infty&\text{altrimenti}
    \end{dcases},
    $$
    i cui autovalori e autofunzioni sono ben noti
    $$
    E^{(1)}_n = \frac{\pi^2\hbar^2}{2mL^2}n(n+2),\quad
    \psi^{(1)}_n(x)=\sqrt{\frac{2}{L}}\sin\left(\frac{(n+1)\pi x}{L}\right),\;
    \text{con }n=0,1,2,3...$$
    Dallo stato fondamentale si può ricavare il superpotenziale corrispondente
    a $V_1$:
    $$W(x) = -\frac{\hbar}{\sqrt{2m}}\frac{\psi'^{(1)}_0(x)}{\psi_0^{(1)}(x)} = 
    -\frac{\hbar\pi}{L\sqrt{2m}}\cot\left(\frac{\pi x}{L}\right).
    $$
\end{frame}
\end{document}
